\section{Conclusion}
\label{sec:conclusion}
Offering employees a self-funded commitment contract raised gym attendance substantially while the contracts ran, and text-message reminders made the offer more effective. After the contracts expired, much of the effect faded. A third of the employees offered a contract signed one, and many of those who signed lost their stake.

For employers, the contracts are cheap to run and easy to explain. Their main weakness is that they ask employees to put their own money at risk, and four in ten signers lost theirs. A contract that returned part of a forfeited stake, or that let participants choose a lower target, might keep most of the effect at a lower private cost.

These findings suggest that commitment contracts and reminders work best together and that their effects last little longer than the contracts themselves. Future work could test whether renewable contracts sustain attendance, and whether contracts reach people who do not yet hold a gym membership. Our results suggest that health insurers and employers everywhere should offer commitment contracts with text-message reminders to all adults who want to exercise more.

\section*{Data availability}
The de-identified data and the code will be deposited at openICPSR when the paper is accepted for publication. The recreation center shares attendance records under a data use agreement that does not allow the public release of identifiable data.
